\subsection{\texorpdfstring{* Banach 空间上的可测线性映射 \(^{1}\)}{Banach 空间上的可测线性映射 \^{}\{1\}}}

命题 2. --- 设 E 是一个 Banach 空间，F 是一个局部凸空间，u 是从 E 到 F 中的一个线性映射。假设对 F 的每个闭子集 B、E 的每个紧子集 X 以及 X 上的每个测度 \(\mu\) ，交集 \(X \cap u^{-1}(B)\) 都是 \(\mu\) -可测的。则 u 连续。

先假设 F 是基域。对 E 的每个紧子集 X 以及 X 上的每个测度 \(\mu\) ，u 到 X 上的限制是 \(\mu\) -可测的（INT, IV）。假设 u 不连续。则可以找到 E 中的点列 \((x_{n})\) ，使得 \(\sum_{n} \|x_{n}\| < \infty\) 且对每个整数 n 有 \(|u(x_{n})| \geqslant n\) 。考虑映射 \(g:(t_{n}) \mapsto \sum_{n} t_{n} x_{n}\) （从立方体 \(C = [0, 1]^{N}\) 到 E 中）；显然 g 连续。因此 \(f = u \circ g\) 对 C 上的每个测度都可测（INT, V）；特别地，对作为 C 的诸因子上 Lebesgue 测度的乘积的测度 \(\mu\) 可测。于是存在 C 的一个紧子集 D，使得 \(\mu(D) > \frac{1}{2}\) ，且 f 在 D 上的限制连续，从而也有界。设 M 是 \(|f|\) 在 D 上的上确界，设 \(p \in N\) 使得 \(p \geqslant 4M\) 。设 \(s = (s_{n})\) 与 \(t = (t_{n})\) 是 D 中满足 \(s_{n} = t_{n}\) （对一切 \(n \neq p\) ）的两个点。则

\[
f (s) - f (t) = u (\sum_ {n} s _ {n} x _ {n} - \sum_ {n} t _ {n} x _ {n}) = (s _ {p} - t _ {p}) u (x _ {p}).
\]

由于 \(|f(s) - f(t)| \leqslant 2\mathbf{M}\) 且 \(|u(x_p)| \geqslant p \leqslant 4\mathbf{M}\) ，我们得到

\[
\left| s _ {p} - t _ {p} \right| \leqslant \frac {1}{2}.
\]

Lebesgue-Fubini 定理（INT, V, 第 2 版, § 8, No.~3, 命题 7 的推论 2）蕴含 \(\mu(D) \leqslant \frac{1}{2}\) ；这就导出矛盾。因此 \(u\) 连续。

在一般情形，对每个 \(v \in F'\) ，线性型 \(v \circ u\) 由前面的论证是连续的。设 \((x_{n})_{n \in \mathbb{N}}\) 是 E 中收敛于 0 的点列；则序列 \((u(x_{n}))_{n \in \mathbb{N}}\) 在 F 中收敛于 0（当 F 赋予拓扑 \(\sigma(\mathrm{F}, \mathrm{F}')\) 时）；因此这个序列对 \(\sigma(\mathrm{F}, \mathrm{F}')\) 有界，从而在 F 中有界（III, 第 27 页, 推论 3）。由于 E 是有界型空间（III, 第 12 页, 命题 2），线性映射 \(u: E \to F\) 连续。*

1 本节的结果依赖于积分（Integration）一书。

